\section{The training pipeline on CPU and Apple MPS}
\label{sec:pipeline}

\subsection{From a game to a deployed checkpoint}
Figure~\ref{fig:pipeline} traces one unit of experience through the pipeline.
Each arrow is a durable hand-off: a SQLite row, a job, a file written by
atomic replacement, or a lease. A crash at any point therefore leaves either
the old state or the new one, and the next supervisor tick resumes from it.

\begin{figure}[htbp]
\centering
\resizebox{\linewidth}{!}{%
\begin{tikzpicture}[x=1cm,y=1cm,
  pstep/.style={box,text width=33mm,minimum height=13mm},
  pgate/.style={warmbox,text width=33mm,minimum height=13mm}]
  \node[pstep] (c) at (0,0) {\textbf{1 Collect} (CPU)\\self-play, sampled\\at $\tau=0.25$};
  \node[pstep] (q) at (4.4,0) {\textbf{2 Queue}\\complete episodes\\and a \code{learn} job};
  \node[pgate] (g) at (8.8,0) {\textbf{3 Learner gates}\\nothing staged; $<2$\\in flight; $\ge256$ steps};
  \node[pgate] (v) at (13.2,0) {\textbf{4 Validate}\\quality checks;\\likelihood $\le10^{-4}$};
  \node[pstep] (p) at (13.2,-2.6) {\textbf{5 PPO} (MPS)\\GAE, 4 epochs,\\KL stop at 0.03};
  \node[pstep] (s) at (8.8,-2.6) {\textbf{6 Save candidate}\\mark consumed;\\freeze inputs; scout};
  \node[pstep] (e) at (4.4,-2.6) {\textbf{7 Evaluate} (CPU)\\incumbent, then\\candidate, paired};
  \node[pgate] (pg) at (0,-2.6) {\textbf{8 Paired gate}\\margin $\lceil nm\rceil$ and\\sign test $p\le0.05$};
  \node[pstep] (r) at (0,-5.4) {\textbf{9 Stage}\\ready slot, SHA-256,\\source generation};
  \node[pgate] (pr) at (4.4,-5.4) {\textbf{10 Promote}\\between games, no\\pending ladder game};
  \node[pstep] (d) at (8.8,-5.4) {\textbf{11 Reload}\\the actor reloads\\by revision};
  \node[box,text width=33mm,minimum height=13mm,draw=muted,fill=white] (rej) at (13.2,-5.4) {\textcolor{muted}{Rejected candidate}\\\textcolor{muted}{kept for audit}};
  \draw[arr] (c)--(q); \draw[arr] (q)--(g); \draw[arr] (g)--(v); \draw[arr] (v)--(p);
  \draw[arr] (p)--(s); \draw[arr] (s)--(e); \draw[arr] (e)--(pg);
  \draw[arr] (pg)-- node[note,left]{pass} (r); \draw[arr] (r)--(pr); \draw[arr] (pr)--(d);
  \draw[darr] (pg.south east) |- node[note,above,pos=.75]{fail} (13.2,-4.0) -- (rej.north);
  \draw[darr] (d.south) -- ++(0,-1.0) -| node[note,below,pos=.25]{the new incumbent collects the next batch} (-2.3,0) -- (c.west);
  \draw[darr] (g.north) to[out=90,in=90,looseness=.6] node[note,above]{defer and retry later} (q.north);
\end{tikzpicture}}
\caption{The lifecycle of experience. Orange boxes are gates that can stop
the flow. Steps 1 and 7 run on the CPU, step 5 on the MPS GPU. Promotion
(step 10) is called by the player between games, never by a worker. Episodes
are marked consumed when the candidate is saved, before it is evaluated, so
``trained'' means \emph{used for a candidate}, not \emph{deployed}.}
\label{fig:pipeline}
\end{figure}

\paragraph{Collection.} The collector snapshots the incumbent and freezes
its inference inputs. It then plays \code{simulation_games} games (default
24) in waves of spawned one-thread processes. Each process drives the
official \code{BattleStream} through \code{simulator.cjs} and observes only
its own player stream. The \code{ladder-v1} curriculum cycles the opponent
through
\[
  (\text{self},\ \text{tactical},\ \text{heuristic},\ \text{self},\ \text{tactical},\ \text{random},\ \text{heuristic},\ \text{tactical}),
\]
with index $(\lfloor i/|\text{teams}|\rfloor+\text{seed})\bmod8$, so one
opponent completes a sweep of the team pool before the next takes over. In
each cycle that is 2/8 self-play against the frozen incumbent, 3/8 tactical,
2/8 heuristic and 1/8 random. Open team sheets are drawn with probability
0.08. Sides alternate as $(i+\lfloor i/|\text{teams}|\rfloor)\bmod2$. Seeds
are 64-bit, split into four 16-bit words. An earlier encoding used only 16
bits and was retired.

\paragraph{Learning.} The learner applies the gates of step~3, the
validation of step~4 and the loss of Section~\ref{sec:learning}. It runs on
the device chosen by \code{backend_for}. In automatic mode MPS is chosen only
if its benchmarked time is at most 0.85 of the CPU time. Any device error
retrains the same batch from the incumbent on the CPU. For the strategic
profiles the learner then refreshes the candidate's embedded
\code{strategy_knowledge} with new ladder postgame counts. It also trains the
public-move scout on the same lane.

\paragraph{Evaluation.} Incumbent and candidate are frozen CPU snapshots.
They play $n$ games each (60 in the deployed configuration; default 20), on
seeds $s_0+100000+i$ derived from the job identity. Actions are sampled at the
checkpoint temperature. Both panels face the same frozen \code{self} snapshot
and the same frozen inputs. Panels resume after interruption. The evaluator
abandons a plan if the source generation, parent revision or candidate hash
has changed.

\subsection{Three lanes on one laptop}
Figure~\ref{fig:lanes} shows how the work is shared out in time. A supervisor
under launchd \code{KeepAlive} wakes every 15\,s. It reloads the autopilot
settings, samples resources, enqueues daily feeds and collection, and starts at
most one child per lane, running \code{python -m ml.worker --lane} with the
lane name (\code{collector}, \code{learner} or \code{evaluator}). Each child takes
its own lock and claims jobs atomically with \code{BEGIN IMMEDIATE}. A job's
lease is its time budget plus 60\,s. A job can be claimed while its attempt
count is below three, and a deferral gives its attempt back, so deferrals do
not exhaust a job.
A child exits after its time budget, and the next tick replaces it.

\begin{figure}[htbp]
\centering
\resizebox{\linewidth}{!}{%
\begin{tikzpicture}[x=1cm,y=1cm]
  \def\W{14.2}
  \foreach \y/\name in {0/{Supervisor\\(CPU)},-1.25/{Collector\\(CPU)},-2.5/{Learner\\(MPS GPU)},-3.75/{Evaluator\\(CPU)},-5.0/{Live runner\\(CPU + Node)}}{
    \draw[lane] (0,\y-0.45) rectangle (\W,\y+0.45);
    \node[lanelabel] at (-0.15,\y) {\name};
  }
  % supervisor ticks
  \foreach \x in {0.3,1.3,...,13.3}{ \fill[accent] (\x,0) circle (1.6pt); }
  \node[note,anchor=west] at (13.45,0) {15\,s};
  % collector
  \fill[pale,draw=ink] (0.3,-1.6) rectangle (4.6,-0.9); \node[font=\scriptsize] at (2.45,-1.25) {rollouts: 24 games, spawned};
  \fill[pale,draw=ink] (8.4,-1.6) rectangle (12.6,-0.9); \node[font=\scriptsize] at (10.5,-1.25) {next collection (new incumbent)};
  % learner
  \fill[warm,draw=seriesB!70!black] (4.8,-2.85) rectangle (6.2,-2.15); \node[font=\scriptsize] at (5.5,-2.5) {PPO};
  \fill[warm,draw=seriesB!70!black] (6.3,-2.85) rectangle (7.0,-2.15); \node[font=\scriptsize] at (6.65,-2.5) {scout};
  \node[note,anchor=north] at (5.9,-2.9) {seconds per batch};
  % evaluator
  \fill[paleink,draw=ink] (7.1,-4.1) rectangle (9.9,-3.4); \node[font=\scriptsize] at (8.5,-3.75) {incumbent panel ($n$)};
  \fill[paleink,draw=ink] (10.0,-4.1) rectangle (12.8,-3.4); \node[font=\scriptsize] at (11.4,-3.75) {candidate panel ($n$)};
  \fill[warm,draw=seriesB!70!black] (12.9,-4.1) rectangle (13.7,-3.4); \node[font=\scriptsize] at (13.3,-3.75) {gate};
  % live runner
  \foreach \a/\b in {0.3/2.6,3.1/5.6,6.1/8.2,8.7/11.3,11.8/13.9}{
    \fill[white,draw=muted] (\a,-5.35) rectangle (\b,-4.65);
  }
  \node[font=\scriptsize,text=muted] at (1.45,-5.0) {game};
  \node[font=\scriptsize,text=muted] at (12.85,-5.0) {game};
  \draw[arr] (4.6,-1.25) to[out=0,in=180] (4.8,-2.5);
  \draw[arr] (7.0,-2.5) to[out=0,in=180] (7.1,-3.75);
  \draw[darr,draw=seriesB] (13.3,-4.1) to[out=-90,in=90] node[note,right,pos=.6,text=seriesB!80!black]{staged} (11.55,-4.62);
  \node[note,text=seriesB!80!black] at (11.55,-5.65) {promoted in the gap between games};
\end{tikzpicture}}
\caption{Schematic timeline of the lanes; durations are not to scale. The
collector and evaluator share one pool of simulator slots. Within the default
budget of four, the evaluator gets half (calibrated locally to three) and the
collector the rest. The learner owns the GPU. The live runner is outside the
pipeline and only ever consumes a staged candidate at a game boundary.}
\label{fig:lanes}
\end{figure}

\subsection{Resource admission}
Every lane re-samples the host before starting work and at minibatch
checkpoints. With $C$ logical cores, a measured CPU load of $u$ percent,
available memory $F$\,GiB, $r$ resident workers and an estimated $m$\,MiB per
worker, the allowed simulator worker count is
\begin{align}
  W_{\mathrm{mem}}&=\max\!\big(0,\lfloor(F-F_{\min})1024/m\rfloor+r\big),\qquad
  W_{\mathrm{cpu}}=\max\!\big(0,\lfloor(u_{\max}-u)C/100\rfloor\big),\\
  W&=\mathbf 1[\text{healthy}]\cdot\max\!\big(0,\min[W_{\max},\,C-C_{\mathrm{res}},\,W_{\mathrm{mem}},\,W_{\mathrm{cpu}}]\big),
\end{align}
where \emph{healthy} means $F\ge F_{\min}$, $u<u_{\max}$, swap-out below its
limit, no critical memory pressure (macOS
\code{kern.memorystatus_vm_pressure_level} flag value~4), and at least 2\,GiB
of free disk. ``Yellow'' pressure (flag value~2) is allowed; only ``red''
blocks. The code
defaults are $C_{\mathrm{res}}=3$, $F_{\min}=2$\,GiB, $W_{\max}=4$,
$u_{\max}=75\%$, $m=400$\,MiB and swap-out 32\,MiB/s. The deployment uses a
modestly higher CPU ceiling and a calibrated swap limit. Optional work also
stops above 90\% of the storage budget (Figure~\ref{fig:storage}).

\begin{proposition}[Headroom invariants]
\label{prop:headroom}
Whenever $W>0$:
\begin{enumerate}
  \item at most $C-C_{\mathrm{res}}$ simulator workers run, so $C_{\mathrm{res}}$
    cores stay free of simulator workers;
  \item if each \emph{new} worker uses at most $m$\,MiB, launching $W-r$ of them
    leaves at least $F_{\min}$ GiB available;
  \item if each worker adds at most $100/C$ percentage points of load (one core),
    the projected load $u+100W/C$ stays at or below $u_{\max}$.
\end{enumerate}
\end{proposition}
\begin{proof}
(1) $W\le C-C_{\mathrm{res}}$. (2) $W-r\le\lfloor(F-F_{\min})1024/m\rfloor$, so
$(W-r)m/1024\le F-F_{\min}$. (3) $W\le\lfloor(u_{\max}-u)C/100\rfloor$ gives
$100W/C\le u_{\max}-u$.
\end{proof}

\intuition{The memory term credits resident workers, whose memory is already
counted in $F$. The CPU term does not credit them, even though their load is
already in $u$. It is therefore conservative and can under-admit while workers
are running. These are cooperative limits enforced by the application, not OS
quotas. Another service can still take the headroom between samples, and the
5-second re-sampling only bounds how long that can go unnoticed.}

\subsection{The MPS lane}
The learner calls
\code{torch.mps.set_per_process_memory_fraction(0.12)}, which caps its Metal
allocations at 12\% of the device's recommended working-set size
\citep{pytorchmps}. Its thread count is limited and it runs at
\code{nice} 10. A GPU duty fraction $d\in(0,1]$ throttles training. After each
optimizer step the learner synchronizes with Metal, measures the step's busy
time $e_k$, and sleeps
\begin{equation}
  s_k=e_k\Big(\frac1d-1\Big).
\end{equation}

\begin{proposition}[Duty cycle]
\label{prop:duty}
Under this rule each step, and therefore any union of steps, has busy fraction
exactly $d$: $\sum_ke_k/\sum_k(e_k+s_k)=d$.
\end{proposition}
\begin{proof}
$e_k+s_k=e_k/d$, so $\sum_k(e_k+s_k)=\sum_ke_k/d$.
\end{proof}
The proposition concerns wall-clock time inside the training loop, measured
after synchronization. It is not a statement about instantaneous GPU
utilization, which the application cannot observe. The default is $d=0.5$.
For these small cached batches, half duty cost several times the full-duty
time: 34.99 vs.\ 7.45\,ms per update in Table~\ref{tab:bench}, which is more
than the factor of two Proposition~\ref{prop:duty} accounts for. The
deployment therefore runs at full duty and relies on the memory cap and the
admission guards (\code{docs/offline-ppo-followup.md}).

\subsection{The tensor cache}
Rebuilding Python feature lists dominated training time, so batches are
converted to tensors once. With $N$ steps and at most $C$ candidates per step,
the resident footprint in float32 is
$N(384\cdot4+C(192\cdot4+1)+8)$ bytes. The extra byte per candidate is its mask
entry and the 8 bytes per step hold the index. The batch is kept on the device
if this is at most $2^{27}$ bytes (128\,MiB). Otherwise a bounded CPU LRU cache
transfers minibatches one at a time.

\begin{example}[When a batch stays resident]
A batch containing a 360-candidate preview has $C=360$, so each step costs
$1536+360\cdot769+8=278{,}384$ bytes. The residency limit is
$\lfloor2^{27}/278{,}384\rfloor=482$ steps. Recent production batches of 404
and 415 steps were resident; a 617-step batch used the CPU cache.
\end{example}

\begin{table}[htbp]
\centering\small
\begin{tabular}{lrr}
\toprule
64-decision batch, up to 360 choices & Before cache (ms/update) & After cache (ms/update) \\
\midrule
CPU & 43.08 & 14.08 \\
MPS, 50\% duty & 121.61 & 34.99 \\
MPS, full duty & 44.06 & 7.45 \\
\bottomrule
\end{tabular}
\caption{Bounded M5 timing windows from \code{docs/m5-learning-orchestration-research.md}.
They were measured while other workloads ran, so they are diagnostic timings
rather than randomized comparisons. A raw kernel benchmark favoured MPS by
only 1.057 vs.\ 1.512\,ms; most of the speedup came from caching.}
\label{tab:bench}
\end{table}

\subsection{Safe reloads of source code}
The \emph{source generation} is a SHA-256 over the entry points,
\code{simulator.cjs}, the lock file and all \code{ml/*.py} files. Rollouts
refuse to run under a mismatched generation. On a change, the supervisor stops
starting children, drains the running ones and exits with code 75 so that
launchd restarts it under the new code. Evaluation plans and staged candidates
record their generation, and promotion refuses a candidate evaluated under
different code. A source change therefore invalidates evaluations; it never
mixes code versions inside one comparison. On 8--9 October, 25 evaluations
were abandoned for this reason and two because the parent had advanced.

\source{\code{ml/pipeline.py} (supervisor, slots), \code{ml/worker.py}
(lanes, collection, learn, evaluate), \code{ml/resources.py}
(\code{ResourcePolicy}, \code{BackgroundGuard}, \code{backend_for}),
\code{ml/batching.py}, \code{ml/reload.py}, \code{ml/simulator.py},
\code{simulator.cjs}, \code{scripts/learning_service.py}.}
